\documentclass[10pt]{article}
\usepackage[a4paper,margin=25mm]{geometry}
\usepackage{amsmath,amssymb,amsthm,mathtools,booktabs,hyperref,xurl,microtype}
\setlength{\emergencystretch}{3em}
\hypersetup{hidelinks}
\newtheorem{theorem}{Theorem}
\newtheorem{lemma}[theorem]{Lemma}
\newcommand{\Q}{\mathbb Q}\newcommand{\Z}{\mathbb Z}
\title{One half-trace square removes an integer block:\\an ordinary-abc quartic in 605 variables}
\author{\(\Z\Q\) Diophantine Project}
\date{24 November 2026}
\begin{document}\maketitle

\begin{abstract}
The source's integral prefix forces its first Pell parameter to be
divisible by three, before any use of integrality of the coordinate
\(c\). The signed shifted-trace Pell theorem then recovers \(c\)
from one rational half-trace square. We replace the entire ordinary-abc
integer certificate for \(c\) by that square, retaining its still-used
\(c^2\) register as an ordinary source coordinate. The complete table
decreases from 645 variables and 564 quadratic residuals to 605 variables
and 528 residuals, with twelve remaining integer predicates. Soundness
and completeness hold on the source-instance projection under the same
retained arithmetic-interface, ordinary rational abc and coding premises.
\end{abstract}

\section{The exact result}

The predecessor is the complete directly scaled common-point table
of~\cite{scaled}. All its arithmetic-interface premises, ordinary
rational abc, and the ineffective sufficiently large odd defining
coefficient remain in force. The externally cited Jones coding theorem
also remains required. The new arithmetic input is the elementary
integrality theorem proved in~\cite{integrality}, not a further
conjectural premise.

\begin{theorem}\label{thm:main}
The complete table supplied here has 605 rational unknowns and 528
quadratic residuals. Its sum of squared residuals has degree exactly
four. Under the same retained premises as the 645/564 predecessor,
its source-instance projection agrees with that predecessor's projection.
It has twelve ordinary-abc integer predicates; integrality of the
source outputs \(c,d\) follows from the retained source equations
and one added half-trace square.
\end{theorem}

The twelve certified inputs are
\[
 b,g,t,m,w,s,h,\gamma,j,i_{\rm half},
       \mathrm{power\_delta},\mathrm{power\_alpha}.
\]
The coordinate \(c\) remains a rational unknown of the polynomial
system. Its separate elliptic integer predicate is removed. No
universality assumption for the historically emended source is added,
and no unconditional undecidability assertion over \(\Q\) is made.

\section{Why the half-trace square certifies integrality}

We record the precise theorem being applied.

\begin{lemma}[Signed shifted-trace Pell integrality]\label{lem:integer}
Suppose \(a\in3\Z\), \(c,d,\zeta\in\Q\), and
\begin{equation}\label{eq:integrality}
 d^2-(a^2-1)c^2=1,\qquad T=d-(a-2)c\in\Z,
 \qquad \zeta^2=2\bigl(ad+(a^2-1)c+1\bigr).
\end{equation}
Then \(c,d\in\Z\). The conclusion includes negative \(a\)
and \(a=0\).
\end{lemma}
\begin{proof}
This is the one-square theorem of~\cite{integrality}. Its denominator
argument is short enough to make the application explicit. Put
\(D=4a-5\), \(\kappa=a-2\), and \(L_a=2a^2-2a-1\).
The first two equations give
\[
 Dc^2-2\kappa Tc-(T^2-1)=0,\qquad
 (Dc-\kappa T)^2=(a^2-1)T^2-D.
\]
The latter rational square root is integral, so \(Dc\in\Z\).
Thus the reduced denominator \(N\) of \(c\) divides \(|D|\).
The half-trace relation becomes
\[
             \zeta^2=2L_ac+2aT+2,
             \qquad 8L_a+3=D(4a+1).
\]
Since \(3\mid a\), the odd integer \(D\) is prime to three,
and \(\gcd(2L_a,N)=1\). The right side has reduced denominator
exactly \(N\). Therefore \(N\) is a square.

The exact signed rational Pell classification in~\cite{integrality}
now excludes every nonintegral orbit. For \(a\ge3\) divisible by
three, its nonintegral orbit has denominator \(4a-5\equiv3\pmod4\).
For \(a=-A\le-3\) with \(3\mid A\), its sole nonintegral
orbit has denominator \(4A+5\equiv2\pmod3\); the possible
additional seed exists only if \(4A+5\) is a square, which this
congruence excludes. At \(a=0\), the only nonintegral possibilities
are \(c=\pm3/5,\pm4/5\), with denominator five. None of these
denominators is a square. Consequently \(c\in\Z\), and
\(d=(a-2)c+T\in\Z\).
\end{proof}

The preliminary signed-source argument~\cite[Section 2]{signed}
uses the certified inputs other than \(c\) to recover integral
\(q,\lambda,S,T_0,r,n,M,a\). In particular the auxiliary power
block first makes \(q\) integral and gives
\(q\equiv\pm1\pmod{b^5-1}\), hence
\(\lambda=(q^4-1)/(b^5-1)\in\Z\). None of this uses the
first-Pell index classification or assumes \(c\in\Z\).
The retained prefix equations then give
\begin{equation}\label{eq:prefix}
 n=q^{16},\quad M=q^{32}=n^2,\quad
 r=S(M-n)+T_0(M-1),\quad a=(wM+1)rsM.
\end{equation}
If \(3\mid q\), then \(3\mid M\). Otherwise
\(n\equiv M\equiv1\pmod3\), so \(3\mid r\). In both
cases \(3\mid a\). This precedes the later power-of-two recovery;
it does not assume the eventual identity for \(q\).

The retained first linear row is
\[
                d=(a-2)c+bw+(4a-5)\gamma.
\]
Thus its shifted trace is integral, using only the twelve retained
integer inputs and the reconstructed \(a\). All hypotheses of
Lemma~\ref{lem:integer} are now available without circularity.

\section{The literal quadratic replacement}

The predecessor already stores \(A_c=ac\). Its \(c\)-certificate
also stores \(C_2=c^2\), used outside that certificate in the source
rows
\[
             d^2-A_c^2+C_2-1=0,\qquad
             E-(2i_{\rm half}+1)C_2=0.
\]
Retain \(C_2\), rename it \texttt{source\_c\_square}, and retain
its defining row \(c^2-C_2=0\), now named as an ordinary source
definition. No new square register is introduced.

Delete every other row of the \(c\)-integer block, including its
private scaled-point equations. Add one rational coordinate \(\zeta\)
and the single quadratic residual
\begin{equation}\label{eq:new}
                     \zeta^2-2ad-2aA_c+2c-2=0.
\end{equation}
The retained equation \(A_c=ac\) makes this exactly the half-trace
square of Lemma~\ref{lem:integer}. Every other source row and every
remaining integer-certificate row is unchanged after the one coordinate
renaming and reindexing.

The full occurrence audit finds exactly 41 deleted rational coordinates,
all private to the discarded rows. The shared reference chart, common
physical point, reciprocal product, contact-square and contact-grid
coordinates remain. The twelve surviving contact offsets keep their
original positions in the same dyadic grid; the omitted contact list
is simply absent. No offset is relabeled and no simultaneous arithmetic
condition is strengthened.

\section{Soundness and completeness on the instance projection}

\begin{proof}[Proof of Theorem~\ref{thm:main}]
For soundness, consider any rational zero of the new table. The
shared rows and the twelve complete remaining integer blocks are
unchanged. The blockwise polynomial reconstructions of~\cite{scaled}
therefore recover each of their original integer certificates. Those
reconstructions require no row from the discarded \(c\) block;
their shared fields are all retained. Thus the twelve listed inputs
are integers.

The preliminary source argument gives the integral prefix
\eqref{eq:prefix}, so \(3\mid a\). The first linear row gives
the integral shifted trace. The retained first norm, \(A_c=ac\),
and \(C_2=c^2\) give its Pell norm. Equation~\eqref{eq:new}
gives its half-trace square. Lemma~\ref{lem:integer} now proves
\(c,d\in\Z\), including every signed and zero-parameter case.
All integer inputs required by the old source soundness theorem have
therefore been recovered. Apply that theorem to obtain the same
source-instance conclusion. This does not require recreating old
\(c\)-certificate auxiliaries while fixing arbitrary surviving
elliptic witnesses.

For completeness, choose the source's canonical positive witnesses.
They have \(B=2r+1\), \(c=Y_B(a)\), \(d=X_B(a)\), and
\[
 ad+(a^2-1)c=X_{2r+2}(a)=2X_{r+1}(a)^2-1.
\]
Set \(\zeta=2X_{r+1}(a)\). It satisfies~\eqref{eq:new}.
The old common-point completeness theorem supplies all thirteen
integer certificates for these positive witnesses; forget the private
\(c\) block and retain its \(c^2\) register as the source square.
Every new row holds. Equivalently, choose the same simultaneous
construction for the twelve remaining inputs. This proves the
claimed source-instance equivalence.

Exactly 37 old residuals and 41 now-unused coordinates are deleted.
One residual and one coordinate are added, so the complete counts are
\[
                    645-41+1=605,\qquad
                    564-37+1=528.
\]
Every residual has degree at most two in the unknowns, with integer
polynomial coefficients in the retained parameters. The unchanged
\(b^2\)-register residual has nonzero quadratic part, and sums of
squares of rational leading quadratic forms cannot cancel. The sum
of squares therefore has degree exactly four.
\end{proof}

\section{Complete-table verification}

The directory supplies the complete sparse table, deterministic compiler,
independent checker, metadata certificate and exact verification receipt.
Run from the repository root:
\begin{center}\small
\texttt{python -B}\enspace\nolinkurl{undecidability/abc_pell_integrality_v108/build_abc_integrality.py}\\
\texttt{python -B}\enspace\nolinkurl{undecidability/abc_pell_integrality_v108/verify_abc_integrality.py}.
\end{center}
The checker imports no compiler. It compares all 526 unchanged rows
under the coordinate renaming, verifies the renamed square row and new
half-trace, audits every occurrence of all 41 removed coordinates,
and checks the complete shared and contact interfaces. It independently
interprets the ten source-prefix rows needed for~\eqref{eq:prefix},
checks all 27 residue cases modulo three, verifies the three symbolic
denominator identities, and checks 32 exact canonical half-trace
fixtures. It also checks coefficient domains, support, degree, counts
and hashes. The quantified signed integrality theorem is proved
in~\cite{integrality}; the finite fixtures do not replace that proof
or certify the retained elliptic arithmetic premises.

\small
\begin{thebibliography}{9}\raggedright
\setlength{\itemsep}{2pt}
\bibitem{scaled} \(\Z\Q\) Diophantine Project,
\emph{Directly scaled private points: a 645-variable ordinary-abc quartic
transfer}, 24 November 2026,
\nolinkurl{undecidability/abc_scaled_point_v106/scaled_private_points.tex};
complete predecessor
\nolinkurl{undecidability/abc_scaled_point_v106/scaled_point_abc_quadrics.json}.
\bibitem{integrality} Project, \emph{A half-trace square replaces the
c-integrality predicate}, 24 November 2026,
\nolinkurl{undecidability/pell_integrality_replacement_v106/pell_integrality_replacement.tex}.
\bibitem{signed} Project, \emph{Both first-Pell signs have the same
coding projection}, 23 November 2026,
\nolinkurl{undecidability/pell_signed_branch_v101/signed_branch_projection.tex}.
\end{thebibliography}
\end{document}
